\subsection{Continuous layered-ray components}

The finite-polygon endpoint discrepancy motivates a separate continuous-curve
calculation. This component retains the original isotropic Fermat functional
and the finite model box. Straight segments represent constant-speed and
vertical legs; circular arcs represent positive affine depth gradients; a
horizontal guide may lie on the physical bottom. The curve used for integration
is distinct from its sampled display polygon. Its branch information does not
set a polygon gradient or production convergence flag.

For a single affine interval, $v(z)=v_0+g(z-z_0)$ with $g>0$, the unbounded
arc travel time is
\begin{equation}
 T=\frac{2}{g}\operatorname{arsinh}
 \left(\frac{gD}{2\sqrt{v_s v_r}}\right),
\end{equation}
where $D$ is the Euclidean endpoint separation and $v_s,v_r$ are the endpoint
speeds. A finite-bottom guide replaces an inadmissibly deep arc when horizontal
range reaches
\begin{equation}
 R_* = \frac{\sqrt{v_b^2-v_s^2}+\sqrt{v_b^2-v_r^2}}{g},
\end{equation}
where $v_b$ is the bottom speed. The implementation includes the homogeneous
and vertical limits. It stores a small angular sweep independently of the
starting angle to avoid cancellation in nearly straight arcs. It does not yet
enumerate general piecewise-layered or three-dimensional branches.

The integration component retains every nodal trilinear basis function $N_j$,
including the lateral model degrees of freedom about an exactly depth-only
base. On each supplied curve it calculates
\begin{equation}
 T=\int\frac{ds}{v},\qquad
 \frac{\partial T}{\partial v_j}=-\int\frac{N_j}{v^2}\,ds,\qquad
 \ell_j=\int N_j\,ds.
\end{equation}
The final quantity supplies the existing DD length mask. Each curve piece is
split at all crossed $x$, $y$, and $z$ grid planes. Sixteen-point quadrature
supplies the result; comparison with eight-point quadrature is a numerical
diagnostic, not an interval error bound. The global depth-only check uses
exact nodal equality. A lateral model change therefore invalidates this
specialization. Unique nodal envelope derivatives also require a unique
relevant ray; supplied tied or incomplete branch status remains unresolved.

Two independent sets of invented controls passed. Across nine homogeneous,
affine, reversed, bottom-guided and nearly homogeneous cases, the largest
integrated-time discrepancy from the separate closed-form calculation is
$5.42\times10^{-19}$ s, and receiver closure is within
$1.27\times10^{-18}$ km. Central differences of the analytic travel time agree
with source-momentum components within $2.63\times10^{-13}$ s/km. Separate
polynomial and trigonometric antiderivatives check every nodal kernel in
homogeneous and affine fixtures. Their largest velocity-partial and nodal-length
differences are $3.80\times10^{-19}$ s$^2$/km and
$1.52\times10^{-18}$ km. The complete 58 and 24 checks, sources and replay
commands are preserved in
\path{data/development/layered_curved_components_001/}.
The documented fresh-directory replay rebuilt both control programs and
reproduced all three scientific result files byte for byte.

Two subsequent tests used reference-selected rays from the authoritative
Loma Prieta geometry: a bottom-guided ray and an interior-turning ray.
Construction completed with six and five pieces, respectively. Both calls
to the unchanged integrator failed because a quadrature node rounded to a
partition endpoint. Neither produced a completed native integral or display.
The largest intended-versus-stored depth-join differences were
$8.67\times10^{-19}$ and $1.30\times10^{-18}$ km. These saved discrepancies
are distinct from the reported integration failure; their causal role needs
the separate partition diagnosis. The complete inputs, identities, selected
references and failures are retained in
\path{data/development/layered_reference_selected_curves_001/}.

These results establish a limited continuous-curve component. No real-data
inversion, updated-model accuracy result, or manuscript-panel promotion follows
from these tests.
